\chapter{Trabalhos Relacionados}

A automação da triagem clínica e da anamnese evoluiu de sistemas de regras para modelos
probabilísticos e, mais recentemente, para arquiteturas baseadas em Transformers.

\section{A Era Pré-Neural: ELIZA e DOCTOR}

O ELIZA, na década de 1960, demonstrou a viabilidade de interação conversacional em saúde
mental com base em padrões linguísticos simples. Apesar da ausência de compreensão
semântica, estabeleceu fundamentos de interação humano-computador em contexto clínico.
\cite{wikipedia2025}

\section{Sistemas de Triagem Probabilística}

Soluções baseadas em redes Bayesianas e grafos causais avançaram o diagnóstico computacional
por inferência probabilística, porém com limitação de interação natural e dependência de entradas
estruturadas \cite{richens2020}.

\section{Agentes Conversacionais em Saúde Mental}

Sistemas como Woebot mostraram que interfaces conversacionais podem apoiar acompanhamento
de saúde mental com impacto mensurável, desde que o desenho conversacional priorize empatia e
clareza \cite{fitzpatrick2017,woebot2025}.

\section{Modelos de Linguagem Biomédicos}

Com BioBERT, observou-se ganho de desempenho em tarefas de mineração de texto clínico,
principalmente NER, quando comparado a modelos generalistas não especializados em domínio
biomédico \cite{lee2020}.

\section{Grandes Modelos de Linguagem na Clínica}

\subsection{Modelos Especialistas}

Modelos como Med-PaLM indicam potencial de alto desempenho em tarefas médicas específicas,
mas exigem validação rigorosa para uso em cenário real \cite{singhal2023}.

\subsection{Modelos Generalistas}

Modelos generalistas apresentam forte capacidade de linguagem, porém trazem desafios de custo,
governança e risco de alucinação factual em aplicações críticas.

\section{Paradigma RAG}

A estratégia Retrieval-Augmented Generation combina geração e recuperação de conhecimento
externo para reduzir alucinações. Esta direção fundamenta a evolução planejada do sistema no
TCC 2 \cite{lewis2020,alura2025,aws2025}.
